\subsubsection{Event-Scheduling Scheme}
Questo approccio si concentra sugli eventi come entità centrali che modificano lo stato del sistema in istanti temporali discreti.

Il simulatore mantiene una lista centrale ordinata chiamata \textbf{Scheduled Event List} (\textbf{SEL}). Ogni elemento della lista è una coppia composta da un tipo di evento e il momento esatto in cui accadrà. Il motore di simulazione funziona estraendo l'evento \textit{più imminente}, portando il tempo globale della simulazione a quel valore ed eseguendo la routine associata a quello specifico evento per aggiornare lo stato del sistema (es. un arrivo o una partenza).

La logica è prettamente legata alle attivazioni temporali e risulta "frammentata". Non esiste una visione globale della vita di un oggetto; ci sono solo blocchi isolati di codice che si attivano quando scatta un determinato clock.

\subsubsection{Process-Oriented Scheduling Scheme}
Questo approccio sposta il focus dagli eventi alle \textbf{entità} (es. i clienti in una coda, i pezzi in una catena di montaggio o i pacchetti in una rete) che attraversano il sistema.

Un \textit{processo} descrive l'intera sequenza di eventi e intervalli temporali che una singola entità sperimenta durante il suo ciclo di vita all'interno del sistema. Ad esempio:
\begin{enumerate}
	\item Arriva l'entità;
	\item Si controlla che il server è libero;
	\item Si riserva la risorsa;
	\item Attesa per la durata del servizio;
	\item Rilascio della risorsa;
\end{enumerate}

Nella simulazione orientata ai processi, ogni entità si comporta come se fosse un oggetto o un thread indipendente con la propria linea di controllo. Il codice viene scritto dal punto di vista dell'entità, e contiene funzioni che possono "sospendere" temporaneamente l'esecuzione del thread (es. funzioni di time delay o attese condizionate in coda prima di ottenere una risorsa).

\subsubsection{Summary schemes}
Molti linguaggi e software di simulazione moderni (come Arena o SIMSCRIPT citati nel testo) offrono un'interfaccia process-oriented perché è molto più intuitiva per il programmatore. Tuttavia, "sotto il cofano" (a livello di kernel software), spesso traducono tutto in un motore event-scheduling con una SEL centrale, perché è computazionalmente molto più efficiente da eseguire.

\begin{table}[H]
	\centering
	\caption{Confronto tra Event-Scheduling Scheme e Process-Oriented Scheme.}
	\label{tab:scheduling_schemes}
	\begin{tabularx}{\textwidth}{l X X}
		\toprule
		\textbf{Caratteristica} & \textbf{Event-Scheduling Scheme} & \textbf{Process-Oriented Scheme} \\
		\midrule
		\textbf{Focus principale} & Gli eventi atomici e isolati (Cosa succede al tempo $t$?) & Le entità e il loro ciclo di vita (Cosa fa questo oggetto nel sistema?) \\
		\addlinespace
		\textbf{Punto di vista del codice} & Frammentato in routine basate sul tipo di evento & Sequenziale e lineare per l'entità (simile alla logica di un thread) \\
		\addlinespace
		\textbf{Gestione del tempo} & Salti espliciti basati sulla \emph{Scheduled Event List} (SEL) & Sospensione del processo tramite funzioni di ritardo (\emph{delay}) \\
		\addlinespace
		\textbf{Flessibilità e utilizzo} & Universale, permette di modellare dinamiche molto complesse e arbitrarie & Ideale per sistemi di accodamento e logistica; richiede molte meno righe di codice \\
		\bottomrule
	\end{tabularx}
\end{table}